\section{Aufbau}

\subsection{NaI-Szintillationsdetektor} \label{sec:nai_aufbau}

Um die Energie und damit die Wellenlänge von $\gamma$-Photonen zu vermessen, können Szintillationsdetektoren genutzt werden. Im durchgeführten Versuch wird dazu ein \textit{NaI:Tl}-Kristall (mit Thallium dotiertes (oder aktiviertes) Natrium-Iodid \cite[559-561]{gammahandbook}) als Szintillator sowie eine angeschlosse Photomultiplier-Röhre (PMT, für \textit{Photo-Multiplier-Tube}) genutzt. Der Aufbau des Detektors (ohne Speicher- und Ausleseelektronik) ist in Abbildung \ref{fig:scintillator_aufbau_skizze} schematisch abgebildet.

Fällt ein $\gamma$-Photon in den Kristall ein, kann es dort via Photoeffekt oder Compton-Streuung (siehe Abschnitt \todo{ein Abschnitt mit WWen erklärt}) ein Elekton der Kristallgitter-Atome ionisieren, welches wiederrum mit mehreren (Anzahl näherungsweise proportional zur Energie des $\gamma$-Photons) gebundenen Gitterelektronen wechselwirkt und diese ins Leiterband des Kristalls anregt. Durch die Dotierung mit Thallium können diese angeregten Elektronen in zusätzliche Energieniveaus (zwischen Leitungs- und Valenzband des NaI-Kristalls) fallen, wobei ein Photon der Frequenz $\nu = \frac{\Delta E}{h}$ (mit $h=\SI{6.626e-37}{\joule\per\hertz}$ als Planck'sches Wirkungsquantum\cite{KonstantenNIST}) emittiert wird. $\Delta E$ ist dabei der Energieunterschied zwischen dem Leitungsbandniveau von NaI und dem Thallium-Niveau in der Bandlücke, also eine Materialkonstante des Szintillators und unabhängig von der Energie des einfallenden $\gamma$-Photons \cite[165-166]{leotechniques}.
Der NaI-Kristall ist transparent gegenüber diesem Szintillationsphoton, welches dann auf die Photokathode zu Beginn der PMT einfallen kann. Dort kann es mit dem Kathodenmaterial wechselwirken, sodass über Photoeffekt ein Photoelektron emittiert wird, welches durch von außen angelegte Hochspannung zur ersten Dynode beschleunigt wird. Dort werden über Stöße des beschleunigten Elektrons mehrere weitere Elektronen ausgeschlagen, welche nun weiter zur nächsten Dynode beschleunigt werden. Über den Verlauf der Röhre und der zunehmend stärkeren Spannung, die an den Dynoden anliegt, steigt der Elektronenstrom sehr stark an. Am hinteren Ende der PMT ist die Anzahl der Elektronen groß genug, um als Stromimpuls gemessen zu werden. Hierbei ist die gemessene Höhe des Stromimpulses annähernd proportional zur Energie des ursprünglichen $\gamma$-Photons, da die Verstärkung der PMT für jedes Szintillationsphoton (bis auf statistische Schwankungen der ausgeschlagenen Elektronen an den Dynoden) identisch ist \cite[563-565]{gammahandbook}.

\jafps{scintillator_aufbau_skizze}{Aufbau des Szintillationsdetektors, Erklärung der einzelnen Bauteile siehe Abschnitt \ref{sec:nai_aufbau}. Abbildung entnommen aus \cite[250]{mcgregorscintillator}}{0.6}

\subsection{HPGe-Halbleiterdetektor}
